
\فصل{مقدمه}

جداسازی کور منابع\LTRfootnote{Blind Source Separation} 
یا 
تحلیل مؤلفه‌های مستقل\LTRfootnote{Independent Component Analysis}
مسئله‌ای  در پردازش سیگنال و یادگیری ماشین است که در اویل دهه‌ی ۸۰ میلادی توسط جوتن\LTRfootnote{Christian Jutten}  مطرح شد
\cite{historicalnotes}. 
در مسئله‌ی جداسازی کور، سیگنال‌های چند منبع مستقل، پس از گذر از یک سیستم با هم مخلوط می‌شوند و این مخلوط‌ها توسط چند گیرنده دریافت می‌گردند. هدف، یافتن منابع اولیه از روی مخلوط‌های آن‌ها و بدون داشتن هیچ اطلاعات اولیه‌ای در مورد منابع یا سیستم مخلوط کننده است. به عنوان مثال،  فرض کنید $n$ نفر در یک اتاق به طور هم‌زمان در حال صحبت هستند.  صحبت این افراد توسط $n$ میکروفن در مکان‌های دلخواهی از این اتاق ذخیره شده است. هدف، جداسازی سیگنال صحبت هر فرد، تنها با مشاهده‌ی سیگنال ذخیره‌شده توسط این میکروفن‌ها، که هر یک مخلوطی از سیگنال صحبت افراد مختلف است، می‌باشد. جداسازی کور منابع می‌تواند در مسائل گوناگون، از جمله جداسازی سیگنال منابع مختلف مغزی در هنگام مطالعه بر روی سیگنال‌های الکتروانسفالوگرام، حذف سایه‌ی پشت یک برگه در هنگام اسکن کردن آن و همچنین مسئله‌ی استنتاج علّی کاربرد داشته باشد.


فرض کنید که  فرآیند‌های تصادفی‌ مستقل
$\s(t) = [s_1(t), s_2(t) \dots, s_n(t)]^\top$،
منابع ما بوده که مشاهده نشده‌اند و فرآیند‌های تصادفی 
$\x(t) = [x_1(t), x_2(t) \dots, x_n(t)]^\top$
مشاهدات ما باشند.   تابع 
$f: \R^n \to \R^n$
تابع مخلوط کننده است و داریم 
$\x(t) = f\Big(\s(t)\Big)$.
در تمام طول این پروژه، فرض شده است که تعداد منابع و مشاهدات برابر $n$ هستند. هدف اصلی مسئله‌ی جداسازی کور منابع، یافتن تابع $f$ و همچنین سیگنال‌های 
$\s(t)$
بر اساس مشاهدات ما از سیگنال‌های 
$\x(t)$
است. این مسئله در حالتی که $f$ یک تابع خطی باشد، به طور کامل حل شده است و شرایط کافی بر روی منابع تا آن‌ها قابل جداسازی باشند به دست آمده است
\cite{ICAbook}.
در حالت غیرخطی، مسئله بسیار دشوارتر است.  در صورتی که فرآیند‌های تصادفی
$\s(t)$
در زمان، وابستگی نداشته باشند، یعنی سفید باشند، می‌توان مثال‌‌های فراوانی از توابع $f$ یافت که این فرآیند‌های مستقل  (که در زمان نیز رابطه‌ای ندارند و 
سفید هستند) را، به فرآیند‌های مستقل تبدیل کند. تا مدت‌ها تصوّر بر این بوده است که اعمال شرط هموار بودن بر تابع 
$f$
جواب این مسئله را یکتا می‌کند، اما در سال 2002، دکتر بابائی زاده مثالی از تابعی هموار ارائه‌ دادند که قادر بود فرآیند‌های سفید مستقل را به فرآیند‌های سفید مستقل تبدیل کند، در حالی که آن‌ها با همدیگر مخلوط ‌شده‌اند. این مثال نشان می‌دهد که در حالتی که فرآیند‌های تصادفی سفید هستند، جداسازی کور منابع غیرخطی مبتنی بر استقلال منابع ممکن نیست
\cite{mbzadehthesis}.
 در سال ۲۰۰۳، در
\cite{harmeling2003kernel}
ایده‌ی استفاده از وابستگی‌های زمانی را برای جداسازی ترکیب‌های غیرخطی فرآیند‌های تصادفی مطرح شد. در این مقاله‌ الگوریتمی برای جداسازی ترکیب‌های غیرخطی، با استفاده از روابط زمانی ارائه شده است که علی‌رغم عملکرد تجربی بسیار امیدوارکننده‌ی آن،  تا کنون اثباتی برای صحت عملکرد آن ارائه نشده‌است.  در سال‌های ۲۰۱۶، ۲۰۱۷ و ۲۰۱۹، تلاش‌های مجددی برای اثبات یک قضیه‌ی جدایی‌پذیری انجام گرفت و  دو ایده‌ی یادگیری بر مبنای جایگشت منابع\LTRfootnote{Permutation Contrastive Learning (PCL)}
 و یادگیری بر مبنای تمایز بین زمان‌ها\LTRfootnote{Time Contrastive Learning (TCL)} 
  مطرح شد
\cite{TCL, PCL, newHyva}.
در این مقالات، قضیه‌هایی جدید برای جدایی‌پذیری ترکیب‌های غیرخطی فرآیند‌های تصادفی ارائه‌ شده‌است. این قضایا، شرایطی بسیار سخت‌گیرانه بر روی فرآیند‌های تصادفی قرار می‌دهد و تا کنون مشخص نیست که این شرایط، تا چه اندازه می‌توانند محدود‌کننده باشند. یکی از اهداف این پروژه، یافتن دسته‌ای از فرآیند‌های تصادفی با شرایط فوق و بررسی میزان محدودیت اعمال شده توسط آن شرایط است. دسته‌ای از ایده‌ها به دنبال یافتن نگاشتی  هستند که تحت آن‌ها، مسئله‌ی جداسازی کور منابع غیرخطی به یک مسئله‌ی خطی تبدیل شود. مقاله‌ی 
\cite{ehsandoust}
یکی از این مقالات است و الگوریتمی ارائه می‌کند که مشتق فرآیند‌های تصادفی منبع را جداسازی می‌کند، اما قادر به اثبات قضیه‌ی جدایی‌پذیری نیست. 

در این پروژه، به بررسی نتایج تئوری و تلاش‌های موجود برای اثبات قضیه‌ی جدایی‌پذیری مخلوط‌های غیرخطی فرآیند‌های تصادفی خطی و غیرخطی می‌پردازیم. نشان می‌دهیم در حالت کلّی، جداسازی مخلوط‌های غیرخطی غیرممکن است و مثال‌هایی از مسائل غیرخطی معرفی می‌کنیم که در آن‌ها مخلوط‌های غیرخطی جدایی‌پذیر نیستند. بعد از مرور الگوریتم‌های کمینه‌سازی اطلاعات متقابل و کاربرد‌های آن‌ها در جداسازی مخلوط‌های خطی، الگوریتمی مشابه برای جداسازی مخلوط‌های غیرخطی ارائه می‌کنیم و عملکرد تجربی این الگوریتم را بررسی می‌نماییم. در ادامه به بررسی معیار‌های استقلال مبتنی بر کرنل‌ها، برای متغیر‌های تصادفی و برای فرآیند‌های تصادفی پرداخته و تلاش می‌کنیم به کمک این ابزار‌ها الگوریتمی جدید و یا قضیه‌ی جدایی‌پذیر‌ی برای مخلوط‌های غیرخطی فرآیند‌های تصادفی ارائه کنیم.

در فصل 
\ref{probdef}
 مسئله‌ی جداسازی کور مخلوط‌های خطی و  غیرخطی را تعریف  و مثال‌های جدایی‌پذیر را بررسی می‌کنیم و ایده‌ی استفاده از روابط زمانی را مطرح می‌نماییم. در فصل 
 \ref{seprability}
قضایای جدایی‌پذیری موجود را در حالت غیرخطی بررسی می‌کنیم. در  فصل
\ref{algs}
مروری بر روش‌های کمینه‌سازی اطلاعات متقابل خواهیم داشت و الگوریتمی برای جداسازی مخلوط‌های غیرخطی ارائه کرده و کیفیت عملکرد آن را می‌سنجیم. در فصل
\ref{kernel}
نیز معیار استقلال
\lr{HSIC}
را معرفی می کنیم که می‌تواند به عنوان جایگزینی برای اطلاعات متقابل در الگوریتم‌های جداسازی مخلوط‌ها به کار برود. در آخر نیز به جمع‌بندی نتایج و بیان ایده‌ها و راه‌های پیش رو برای  پروژه‌ی کارشناسی ۲ می‌پردازیم.


